% texte.tex — f167r, vv. 380–422 (conversion initiale des travaux antérieurs)
\vers[n=380]{ΜΙ.}{Mika}{ἐγώ.}{Moi.}
\vers[suite]{ΚΟ.}{Coryphée}{περίθου νυν τόνδε πρῶτον πρὶν λέγειν.}{Alors commence par ceindre cette couronne, avant que parler. Paix !}
\vers[n=381, metre={ia⁴}]{}{}{σίγα σιώπα, πρόσεχε τὸν νοῦν· χρέμπτεται γὰρ ἤδη}{Silence! Fais attention ! Car voici qu'elle tousse}
\vers[n=382]{}{}{ὅπερ ποιοῦσ᾽ οἱ ῥήτορες. μακρὰν ἔοικε λέξειν.}{comme font les orateurs. Elle paraît en avoir long à dire.}
\vers[n=383]{ΜΙ.}{Mika}{φιλοτιμίᾳ μὲν οὐδεμιᾷ μὰ τὼ θεὼ}{Aucun motif d'ambition, non par les deux déesses,}
\vers[n=384]{}{}{λέξουσ᾽ ἀνέστην ὦ γυναῖκες· ἀλλὰ γὰρ}{ne m'a fait lever pour prendre la parole, ô femmes. Mais}
\vers[n=385]{}{}{βαρέως φέρω τάλαινα πολὺν ἤδη χρόνον}{je suis excédée, depuis longtemps, hélas !}
\vers[n=386]{}{}{προπηλακιζομένας ὁρῶς ἡμᾶς ὑπὸ}{de nous voir traînées dans la boue par}
\vers[n=387]{}{}{Εὐριπίδου τοῦ τῆς λαχανοπωλητρίας}{Euripide, le fils de la vendeuse d'herbes,}
\vers[n=388]{}{}{καὶ πολλὰ καὶ παντοῖ᾽ ἀκουούσας κακά.}{et constamment en butte à des injures de tout genre.}
\vers[n=389]{}{}{τί γὰρ οὗτος ἡμᾶς οὐκ ἐπισμῇ τῶν κακῶν;}{Car quelle est l'injure dont il ne nous salisse ?}
\vers[n=390]{}{}{ποῦ δ᾽ οὐχὶ διαβέβληχ᾽, ὅπουπερ ἔμβραχυ}{Où s'est-il fait faute de nous calomnier, pour peu simplement}
\vers[n=391]{}{}{εἰσὶν θεαταὶ καὶ τραγῳδοὶ καὶ χοροί,}{qu'il y eût des spectateurs, des tragédiens et des choeurs ?}
\vers[n=392]{}{}{τὰς μοιχοτρόπους, τὰς ἀνδρεραστρίας καλῶν,}{nous traitant de galantes, passionnées d'hommes,}
\vers[n=393]{}{}{τὰς οἰνοπίπας, τὰς προδότιδας, τὰς λάλους,}{sujettes au vin, traitresses, bavardes,}
\vers[n=394]{}{}{τὰς οὐδὲν ὑγιές, τὰς μέγ᾽ ἀνδράσιν κακόν·}{rien qui vaille, grand fléau des hommes.}
\vers[n=395]{}{}{ὥστ᾽ εὐθὺς εἰσιόντες ἀπὸ τῶν ἰκρίων}{Si bien qu'à peine rentrés du théâtre,}
\vers[n=396]{}{}{ὑποβλέπουσ᾽ ἡμᾶς σκοποῦνταί τ᾽ εὐθέως}{ils nous soupçonnent, et guettent aussitôt}
\vers[n=397]{}{}{μὴ μοιχὸς ἔνδον ᾖ τις ἀποκεκρυμμένος.}{s'il n'est pas quelque amant caché.}
\vers[n=398]{}{}{δρᾶσαι δ᾽ ἔθ᾽ ἡμῖν οὐδὲν ὧνπερ καὶ πρὸ τοῦ}{Et plus rien ne nous est permis de ce que nous faisions autrefois,}
\vers[n=399]{}{}{ἔξεστι· τοιαῦθ᾽ οὗτος ἐδίδαξεν κακὰ}{tant il a enseigné de vilaines choses}
\vers[n=400]{}{}{τοὺς ἄνδρας ἡμῶν· ὥστ᾽ ἐάν τις ξυμπλέκῃ}{à nos époux. Une femme tresse-t-elle}
\vers[n=401]{}{}{γυνὴ στέφανον, ἐρᾶν δοκεῖ· κἂν ἐκβάλῃ}{une couronne, on la croit amoureuse ; laisse-t-elle tomber}
\vers[n=402]{}{}{σκεῦός τι κατὰ τὴν οἰκίαν πλανωμένη,}{un ustensile en allant et venant par la maison,}
\vers[n=403]{}{}{ἁνὴρ ἐρωτᾷ· `τῷ κατέαγεν ἡ χύτρα;}{le mari demande : « Pour qui brise-t-on la marmite ?}
\vers[n=404]{}{}{οὐκ ἔσθ᾽ ὅπως οὐ τῷ Κορινθίῳ ξένῳ.'}{Nul doute que ce ne soit pour l'étranger de Corinthe ».}
\vers[n=405]{}{}{κάμνει κόρη τις, εὐθὺς ἁδελφὸς λέγει·}{Une fille est-elle souffrante, aussitôt son frère de dire :}
\vers[n=406]{}{}{`τὸ χρῶμα τοῦτό μ᾽ οὐκ ἀρέσκει τῆς κόρης.'}{« Ce teint ne me plaît guère chez la fillette ».}
\vers[n=407]{}{}{εἶἑν, γυνή τις ὑποβαλέσθαι βούλεται}{Mais passons. Une femme privée d'enfants veut}
\vers[n=408]{}{}{ἀποροῦσα παίδων, οὐδὲ τοῦτ᾽ ἔστιν λαθεῖν.}{en supposer un ; même cela ne se peut celer,}
\vers[n=409]{}{}{ἅνδρες γὰρ ἤδη παρακάθηνται πλησίον·}{car les maris se tiennent tout proche.}
\vers[n=410]{}{}{πρὸς τοὺς γέροντάς θ᾽ οἳ πρὸ τοῦ τὰς μείρακας}{Et auprès des vieillards, qui auparavant épousaient les jeunesses,}
\vers[n=411]{}{}{ἤγοντο, διαβέβληκεν, ὥστ᾽ οὐδεὶς γέρων}{il nous a calomniées, tellement que pas un homme d'âge}
\vers[n=412]{}{}{γαμεῖν ἐθέλει γυναῖκα διὰ τοὔπος τοδὶ}{ne consent à prendre femme, en raison de ce vers :}
\vers[n=413]{}{}{`δέσποινα γὰρ γέροντι νυμφίῳ γυνή'.}{« Pour un époux âgé, la femme est un tyran ».}
\vers[n=414]{}{}{εἶτα διὰ τοῦτον ταῖς γυναικωνίτισιν}{En outre, à cause de lui, sur les appartements de femmes,}
\vers[n=415]{}{}{σφραγῖδας ἐπιβάλλουσιν ἤδη καὶ μοχλοὺς}{ils appliquent des cachets et des barres}
\vers[n=416]{}{}{τηροῦντες ἡμᾶς, καὶ προσέτι Μολοττικοὺς}{pour nous garder, et par là-dessus ils entretiennent des chiens molosses,}
\vers[n=417]{}{}{τρέφουσι μορμολυκεῖα τοῖς μοιχοῖς κύνας.}{épouvantails pour les galants.}
\vers[n=418]{}{}{καὶ ταῦτα μὲν ξυγγνώσθ᾽. ἃ δ᾽ ἦν ἡμῖν πρὸ τοῦ}{Encore lui passerait-on cela : Mais ce qu'on nous permettait autrefois,}
\vers[n=419]{}{}{αὐταῖς ταμιεῦσαι καὶ προαιρούσαις λαβεῖν}{de diriger nous-mêmes le ménage et de prendre dans le cellier}
\vers[n=420]{}{}{ἄλφιτον ἔλαιον οἶνον, οὐδὲ ταῦτ᾽ ἔτι}{de la farine, de l'huile, du vin, cela même ne nous est plus}
\vers[n=421]{}{}{ἔξεστιν. οἱ γὰρ ἄνδρες ἤδη κλῄδια}{possible : Car les maris portent sur eux de petites clefs}
\vers[n=422]{}{}{αὐτοὶ φοροῦσι κρυπτὰ κακοηθέστατα}{secrètes, les plus malicieuses du monde,}
